% !TeX root = ../papers/nc-p3.tex
\section{階数15の三次曲線束の全点と添字保持同型曲線}
\label{sec:rank-fifteen-cubic-pencils}
\begin{theorem}\label{thm:rank-fifteen-cubic-pencils}
正の固定曲線の通常な階数15窓の座標を $(x,y)$ とする。
3の上の付値で $X=\bar x\ne0,1$、$Y=\bar y\ne0$ が定義され、
$B=Y/X^3\in\mathbb F_9^\times$ なら、窓は標準Koszul左右AS正則である。
各始点は共通の四つの中心四次式の生成する多項式環上で自由階数256である。
従って $\lambda\ne0$ が3の上に剰余次数2以下の素点を持つ数体に属するとき、
曲線 $y=\lambda x^3$ の通常な全点でこの結論が成立する。$\lambda$ は任意の非零有理数でもよい。
\end{theorem}
\begin{proof}
$X\in\mathbb F_9$ と $Y/(1-X)^2\in\mathbb F_9$ の場合は既知なので除く。
$c^2=X$, $(v^3-v)^2=B$, $K=1/v$, $H=c(K-v)$ とする。
固定した $B$ では第\ref{sec:rank-fifteen-bivariate-word-paths}節の語の偶奇降下により、
各 $v$ の表示を純粋な一変数有理関数体 $\mathbb F_3(v)(c)$ で計算できる。
平方根の積を含む二次函数体は不要である。
五本の中心語の第1,2,3,5列を用いる。任意の幾何 $c$ を扱うため、消去で
逆数を用いた全ての分子と入力分母を平方自由な多項式 $D_v(c)$ に記録する。
各表示の始点0の四中心商は長さ1から13まで
\[
             4,10,20,31,40,44,40,31,20,10,4,1,0
\]
となる。これは $D_v(c)\ne0$ の各点で有効である。

計算の被覆を明記する。$i^2=-1$ として五つの代表 $B$ を取り、
非素体の三値の三乗共役も含める。二つの表示は $B=-1$ で $v,v+1$、
それ以外で $v,-v$ とする。さらに $\mathbb F_3(B)$ を固定する係数のFrobenius共役を全て使う。
$c$ の符号は変えない。全表示の除外多項式の最大公約元から
$X\in\mathbb F_9$ または $Y/(1-X)^2\in\mathbb F_9$ の根を除く。
後者には $N=Bc^6$, $D=(1-c^2)^2$ の $N^9D-ND^9$ を用いる。
非素体なら残りとその三乗共役の最小公倍式を取り、素体上の多項式にする。
\[
\begin{array}{c|rr|r|l}
 B &\deg D_v&\deg D_{v'}&\text{残余次数}&
   \text{既約因子の次数（括弧内は本数）}\\ \hline
 1&3418&3824&18&3(6)\\
 -1&3460&3963&38&4(8),\,6(1)\\
 -1+i&4058&4058&60&4(11),\,8(2)\\
 i&3056&3075&0&\\
 1+i&4097&3741&40&4(10)
\end{array}
\]
有限例外の各既約因子について、分解体で全根 $c$ を
$\mathbb F_3(B)$ のFrobenius軌道に分ける。
各代表で $X=c^2$, $Y=BX^3$ を戻し、必要な平方根が全て分解するまで体を延長して
元の四窓を組み立てる。既知の場合を除く全代表で、四次中心と未切断の五次までの道を再計算し、
四始点とも四中心商が十三次の消滅まで標準次元であることを確認する。
\path{computations/verify_rank_fifteen_cubic_pencils.py} は入力から中心式も再生成する。
五行の主元、最大公約元、全例外の直接計算は有限体での厳密演算である。

$c$ を固定した表示の変更は第\ref{sec:rank-fifteen-oriented-eighth-cut}節の
$\bar p=c^3$ と $q$ の符号選択により、指定した添字保持窓に持ち上がる。
同じ $(x,y)$ の各移動にもこの証明を適用すると、始点ごとに十三次までの次元上界を
与える四中心が取れる。中心四平面のGrassmann多様体上でこれら四つの非空開集合を交差させ、
共通の四中心を得る。命題\ref{prop:klein-central-strata}により、上界は全次数の自由性と
左右AS正則性を与える。数体の主張では所定の付値を係数体へ延長し、$x$ の負付値、
$X=0,1$、および $x$ が単元で $y$ が非単元の場合を既証の全水準の定理で除く。
残る単元域では $B=\bar\lambda$ である。
\end{proof}

\begin{corollary}
上の三次曲線束の添字保持同型不変量の曲線は、$\lambda\ne0,1/27$ のとき種数2の
滑らかな完備曲線の開部分である。通常な全点は最小周期2を持つ。
一位置移動で割った曲線の種数は1である。$\lambda=1/27$ では被約曲線の正規化は有理で、
$(x,p,q)=(3,0,0),(3/4,0,0)$ が周期1の点である。
\end{corollary}
\begin{proof}
定理\ref{thm:rank-fifteen-global-invariants}の座標では $p^2=x^3-27y$、
$q^2=(1-x)^2-4y$ であり、一位置移動は $(p,q)$ の両符号を反転する。
$a=1-27\lambda\ne0$ とすると、$t=p/x$ で $x=t^2/a$, $p=t^3/a$ と復元できる。
$q^2=(1-t^2/a)^2-4\lambda t^6/a^3$ の右辺は平方自由である。
実際、$x$ に関する三次式の判別式は $16\lambda(1-27\lambda)$ で、定数項は1である。
六単根を持つ二重被覆の種数は2となる。移動の不変量 $x,tq$ は
$(tq)^2=ax((1-x)^2-4\lambda x^3)$ の滑らかな種数1曲線を与える。
通常な領域では $x\ne0$ なので $p\ne0$ である。特殊値では $p^2=0$ となり、
被約曲線は $q^2=-(x-3)^2(4x-3)/27$ である。
$q/(x-3)$ が正規化の有理パラメータとなり、$q=0$ の二点とその他で移動の周期が分かれる。
\end{proof}
